% Source: https://artofproblemsolving.com/community/c3h2576098_flaw_in_geogebra
% Stable identity: aops-topic-2576098-q01
No genuine (nondegenerate) triangle exists.

Suppose $AB=AC=6$ and $BC=12$. Then $A$ lies on the circle of radius $6$ about $B$ and on the circle of radius $6$ about $C$. Since $BC=6+6$, these circles touch at exactly one point, the midpoint $M$ of $BC$. Indeed, if a common point $P$ were not on line $BC$, the triangle inequality in triangle $PBC$ would give $BC<PB+PC=12$, a contradiction; and the only point of line $BC$ at distance $6$ from both $B$ and $C$ is $M$.

\begin{asy}
usepackage("amsmath");
size(6cm);
defaultpen(fontsize(10pt));
pair B=(0,0),C=(12,0),M=(6,0);
draw(circle(B,6),gray);draw(circle(C,6),gray);
draw(B--C,linewidth(0.9));dot(B);dot(C);dot(M);
draw(B--(6*dir(60)),dashed);label("$6$",3*dir(60),NW);
draw(C--(C+6*dir(120)),dashed);label("$6$",C+3*dir(120),NE);
label("$B$",B,W);label("$C$",C,E);label("$M$",M,S);label("$12$",(9,0),S);
\end{asy}

So the only possible position of $A$ is $M$: the three ``vertices'' are collinear and the figure has area $0$. This is a degenerate triangle, not a real one. In general, positive lengths $a\le b\le c$ form a nondegenerate triangle exactly when $a+b>c$; here $6+6=12$ is not greater than $12$. So the answer is $\boxed{\text{no}}$.
